% SAFE-RAG Research Grant Proposal -- revised July 2026
% Rebuilt to match the original Word/markdown-style layout:
% A4, 2 cm margins, Cambria-metric serif (Caladea), 11 pt body,
% blue numbered headings with a dark-blue rule, running header, page footer.
% Compile with: xelatex SAFE_RAG_grant_proposal_revised.tex  (twice, for LastPage)

\documentclass[11pt,a4paper]{article}

\usepackage{fontspec}
\usepackage[a4paper,left=2cm,right=2cm,top=2.06cm,bottom=2.2cm,
            headheight=14pt,headsep=0.30cm,footskip=1.05cm]{geometry}
\usepackage{xcolor}
\usepackage{colortbl}
\usepackage{array}
\usepackage{enumitem}
\usepackage{titlesec}
\usepackage{fancyhdr}
\usepackage{lastpage}
\usepackage{microtype}
\usepackage{ragged2e}
\usepackage{needspace}

% ---------------------------------------------------------------- fonts
% Caladea is metrically compatible with Cambria, the font of the source file.
\setmainfont{Caladea}[
  Path = /usr/share/fonts/truetype/crosextra/,
  UprightFont = Caladea-Regular.ttf,
  BoldFont = Caladea-Bold.ttf,
  ItalicFont = Caladea-Italic.ttf,
  BoldItalicFont = Caladea-BoldItalic.ttf ]
\newfontfamily\symbolfont{DejaVu Sans}

% ---------------------------------------------------------------- colours
\definecolor{headblue}{RGB}{46,116,181}   % heading text  (Word Heading 1)
\definecolor{ruleblue}{RGB}{31,56,100}    % heading rule / title / table head
\definecolor{cellgrey}{RGB}{238,242,247}  % metadata table label column
\definecolor{hdrgrey}{gray}{0.8}          % running-header rule

% ---------------------------------------------------------------- spacing
\linespread{1.00}
\frenchspacing
\setlength{\parindent}{0pt}
\setlength{\parskip}{10pt}
\setlength{\arrayrulewidth}{0.5pt}
\setlength{\emergencystretch}{2em}

\setlist[itemize]{leftmargin=36pt,labelwidth=18pt,labelsep=0pt,align=left,
                  itemsep=5pt,parsep=0pt,topsep=6pt,partopsep=0pt,
                  label={\symbolfont\fontsize{7.5}{7.5}\selectfont\char"25CF}}

% ---------------------------------------------------------------- headings
\titleformat{\section}[hang]
  {\color{headblue}\fontsize{13}{15}\bfseries}{\thesection.}{0.35em}{}
  [{\vspace{2.5pt}\color{ruleblue}\titlerule[0.75pt]}]
\titlespacing*{\section}{0pt}{15pt}{4pt}

\titleformat{\subsection}[hang]
  {\fontsize{11.5}{14}\bfseries\itshape}{\thesubsection}{0.5em}{}
\titlespacing*{\subsection}{0pt}{13pt}{5pt}

% ---------------------------------------------------------------- head/foot
\pagestyle{fancy}
\fancyhf{}
\fancyhead[R]{\fontsize{8}{10}\selectfont SAFE-RAG Research Grant Proposal}
\fancyfoot[C]{\fontsize{8}{10}\selectfont Page \thepage\ of \pageref{LastPage}}
\renewcommand{\headrule}{\color{hdrgrey}\hrule height 0.5pt}
\renewcommand{\footrulewidth}{0pt}

% ---------------------------------------------------------------- reference list
\newcommand{\refitem}[2]{%
  \par\noindent\hangindent=19pt\hangafter=1
  \makebox[19pt][l]{[#1]}#2\par}

\begin{document}
\fontsize{11}{15}\selectfont

% ================================================================ title block
{\centering
\vspace*{-2pt}
{\color{ruleblue}\fontsize{10}{12}\bfseries RESEARCH GRANT PROPOSAL}\par
\vspace{6pt}
{\color{ruleblue}\fontsize{16}{19}\bfseries SAFE-RAG: A Schema-Grounded and
Faithfulness-Aware Self-Correcting Retrieval-Augmented Generation Framework for
Trustworthy Financial Compliance Question Answering\par}
\par}

\vspace{12pt}

% ================================================================ metadata table
{\setlength{\parskip}{0pt}\fontsize{10}{13}\selectfont
\renewcommand{\arraystretch}{1.45}
\setlength{\tabcolsep}{6pt}
\begin{tabular}{|>{\columncolor{cellgrey}\bfseries}p{108pt}|p{328pt}|}
\hline
Principal Investigator & Adnan Mashrur Sadad \\ \hline
Supervisor & Dr. Siti Nur Khadijah Aishah Ibrahim \\ \hline
Host Institution & Malaysia-Japan International Institute of Technology (MJIIT),
Universiti Teknologi Malaysia, Kuala Lumpur, Malaysia \\ \hline
Proposed Duration & 12 months \\ \hline
Submission Date & July 2026 \\ \hline
\end{tabular}\par}

\vspace{4pt}

% ================================================================ 1
\section{Executive Summary}

Retrieval-augmented generation (RAG) has become the default approach for question
answering over financial filings, regulatory rulebooks, and audit documentation,
because it grounds model output in retrieved evidence rather than relying on
parametric memory alone. In financial compliance settings, however, a free-text
answer is rarely sufficient: a compliance officer or downstream reporting system
needs a structured judgment, for example a clause identifier, an obligation status,
and a supporting citation, that can be parsed without further human interpretation.

A scoping review of 35 papers (2021-2026) conducted for this proposal, spanning
schema-constrained generation, RAG faithfulness evaluation, agentic self-correction,
and financial document question answering, found that no existing system jointly
performs schema validation, evidence-faithfulness checking, and typed corrective
routing in a financial or regulatory compliance setting. Thirty of those papers were
identified in the original May-July 2026 search; a further five were added by a
July 2026 update sweep of arXiv listings, and three of the five materially narrow the
gap claimed by an earlier draft of this proposal, as recorded in Sections 3, 4 and 8.1.
Systems that guarantee schema-valid output make no claim about whether the values
inside that output are actually supported by evidence; systems that check faithfulness
are evaluated almost exclusively on free text and cannot detect deceptive grounding, a
documented failure mode in which a response is faithful to some retrieved passage but
wrong because that passage concerns the wrong entity or clause.

This project proposes SAFE-RAG, a framework that classifies each generated answer as
schema-invalid, faithful-but-misattributed, or valid, and routes each condition to a
distinct corrective action: local regeneration, entity-constrained re-retrieval, or
abstention. The typed-routing paradigm the design adopts was first demonstrated in a
tool-calling agent setting, where it raised recovery rates from 32.76\% to 89.68\%
[11], and was carried into retrieval-augmented generation during 2026 [19], [20], [21].
That typed routing outperforms undifferentiated retry is therefore taken here as an
established premise rather than claimed as a finding of this project. What remains
untested, and what this proposal claims, is the pair of conditions being routed over:
every published RAG diagnosis vocabulary is trajectory-oriented and can express neither
schema invalidity against a target output schema nor entity-level misattribution, and
none of this work has been evaluated on financial or regulatory content. SAFE-RAG will
be evaluated against pre-registered, falsifiable hypotheses on the ObliQA regulatory
corpus and FinanceBench-derived financial filings, under a correction budget matched
across every compared condition. The project is scoped to deliver a publishable
contribution within 12 months, with an initial deployment context aligned to
Malaysia's regulatory-technology development priorities. [18]

% ================================================================ 2
\section{Background}

Financial institutions increasingly rely on large language models to support regulatory
compliance and document analysis. Retrieval-augmented generation improves factual
grounding over parametric generation alone and has become the default architecture for
question answering over structured financial and regulatory text, building on an
established benchmark landscape of numerical and tabular reasoning datasets such as
FinQA [1], TAT-QA [2], ConvFinQA [3], and MultiHiertt [4], and alongside
finance-specialized language models such as BloombergGPT [5].

Two distinct failure surfaces nonetheless remain unresolved in this setting. First,
generated output can fail to conform to the structured schema a downstream compliance
system expects. Second, output can conform to that schema while still being unsupported
by, or misattributed to, the wrong part of the retrieved evidence. Existing research
addresses these two surfaces in separate literatures: work on schema-constrained
decoding such as XGrammar [6] guarantees syntactic validity but makes no claim about
evidential correctness, while work on faithfulness and hallucination detection such as
RAGAS [7], AlignScore [8], and FRANQ [9] checks whether a claim is entailed by retrieved
context but does not verify that the evidence concerns the correct entity, clause, or
category.

% ================================================================ 3
\section{Problem Statement}

Current RAG systems deployed in financial and regulatory compliance settings cannot
reliably distinguish between two failure modes: \textbf{structural failure}, where the
output does not parse against the required schema, and \textbf{attribution failure},
where the output is schema-valid and passes standard faithfulness checks, yet cites
evidence about the wrong entity, clause, or category. This second failure mode, recently
termed \textit{deceptive grounding} [10], is invisible to standard faithfulness metrics
because those metrics ask only whether a claim is supported by some evidence, not
whether that evidence is about the correct subject.

Compounding this, systems that do detect a failure typically respond with the same
generic retry action regardless of what actually went wrong. That typing a failure
before correcting it outperforms generic retry is, however, no longer an open question.
It was first demonstrated on a tool-calling agent benchmark, where typed,
failure-specific correction achieved 89.68\% recovery against 32.76\% for generic retry
[11], and during 2026 the same paradigm was carried into retrieval-augmented generation.
One system diagnoses the failure type and the earliest failure point in an agentic RAG
trajectory and then applies a repair conditioned on that diagnosis, reporting higher
answer accuracy at lower token cost than rerun-style recovery [19]; a second reaches a
comparable result by routing on a probed hidden-state failure state [20]; and a third
evaluates diagnose-and-repair under an explicitly fixed correction budget [21]. This
proposal therefore adopts typed routing as a premise of its design rather than
presenting it as a finding.

What remains unresolved is the vocabulary in which failures are diagnosed. Every
existing RAG diagnosis vocabulary is trajectory-oriented, consisting of categories such
as retrieval insufficiency, reasoning error, and malformed tool call. None diagnoses
schema validity against a target output schema, and none performs entity-level
attribution checking. The second omission is not incidental. A deceptive-grounding
failure is, by construction, a trajectory in which retrieval succeeded, reasoning was
sound, and every claim is entailed by retrieved text, so it is invisible to a
sufficiency-based diagnosis model for exactly the reason it is invisible to a
faithfulness metric [10]. A diagnosis vocabulary built for trajectory repair cannot
express it. The closest available vocabulary is a peer-reviewed taxonomy of the error
types that occur in realistic RAG systems, released with an annotated error dataset and
a taxonomy-aligned automatic evaluator [22], but that work stops at classification and
offline evaluation and does not drive an inference-time controller. No system identified
in this review applies typed failure classification to the specific combination of
schema-validity and evidence-faithfulness failures that arises in structured-output RAG
for financial and regulatory compliance.

% ================================================================ 4
\section{Research Gap}

A capability coverage analysis of the 35 reviewed studies, scored against five
dimensions, schema-constrained structured output (C1), faithfulness and
evidence-grounding (C2), agentic retrieval control (C3), typed error routing (C4), and
financial or regulatory compliance domain application (C5), returns two findings. The
first is a correction. C4 is no longer empty outside the tool-calling domain: three of
the five studies added by the July 2026 update sweep score full coverage on typed error
routing within retrieval-augmented generation [19], [20], [21]. An earlier draft of this
proposal claimed that C4 was unoccupied in the RAG setting, and that claim is withdrawn.
The second finding is unaffected by those additions: no reviewed study scores full
coverage on C1, C2 and C4 simultaneously. The one reviewed system that couples schema
structure with evidence attribution does so in the regulatory domain but has no typed
routing mechanism [23]; the three RAG typed-routing systems implement neither schema
validation nor entity-level attribution checking; and the original typed-routing system
attains C4 outside RAG entirely [11]. That empty three-way conjunction, rather than typed
routing in general, is what this proposal claims.

Three concrete gaps follow from this pattern, and they are complementary rather than
independent:

\begin{itemize}
\item \textbf{Gap 1:} schema validity and faithfulness are addressed by largely disjoint
literatures that rarely cite each other and have not been combined into a single
validation step that checks both properties of the same output. The single exception
among the reviewed studies enforces a per-rule schema by injecting it at generation time
in order to make attribution checkable in regulatory compliance question answering [23];
because that constraint is applied uniformly to every query, the two properties are never
diagnosed separately as conditions of the same generated output.

\item \textbf{Gap 2:} typed error routing outperforms generic retry, and this has now
been shown inside retrieval-augmented generation as well as for tool-calling agents [11],
[19], [20], [21]. What has not been shown is routing over the pair of conditions targeted
here. Every published diagnosis vocabulary types failures at the level of the agent
trajectory; none diagnoses schema invalidity against a target output schema, none checks
whether cited evidence concerns the entity, clause, or category named in the query, and
none of this work is evaluated on financial or regulatory content.

\item \textbf{Gap 3:} financial RAG systems invest heavily in retrieval engineering,
fine-tuned retrievers, hybrid search, reranking, but remain single-pass at inference,
with no per-query, quality-triggered repair loop. The closest competitor intervenes on
attribution in the regulatory domain by constraining generation uniformly rather than by
repairing diagnosed failures [23], so it never has to decide which condition failed: it
has no schema-invalidity branch, no re-retrieval branch, no abstention branch, and it
reports no comparison against an undifferentiated retry baseline.
\end{itemize}

Together, these gaps motivate a framework that diagnoses schema validity and evidence
faithfulness as two separately detected conditions, transplants typed routing onto that
specific pair of conditions, and targets financial compliance question answering
specifically, where institutional demand for structured, auditable output is already
documented [15], [16] and existing benchmarks remove the need to construct a new corpus
from scratch.

% ================================================================ 5
\section{Research Questions}

\textbf{RQ1.} Can a schema validator and an evidence-attribution checker operate as two
independently diagnosable conditions over structured RAG output in the financial
compliance domain, such that each condition can be detected without relying on the other?

\textbf{RQ2.} Given that typed routing is already established to outperform
undifferentiated retry over trajectory-level failure taxonomies [11], [19], [20], does
the specific two-condition typing proposed here, schema invalidity routed to local
regeneration and evidence misattribution routed to entity-constrained re-retrieval,
outperform an undifferentiated single-retry baseline that uses the same validator and
checker components, under a matched correction budget?

\textbf{RQ3.} Does the proposed evidence-attribution check reduce
deceptive-grounding-style errors, defined as schema-valid, faithfulness-metric-passing
answers that cite evidence concerning the wrong entity, clause, or category, relative to
standard faithfulness checking alone?

% ================================================================ 6
\section{Research Objectives}

\textbf{RO1.} Complete a PRISMA 2020-compliant systematic literature review that extends
the 35-paper scoping review already conducted for this proposal, and finalize the schema
design and adversarial test subset.

\textbf{RO2.} Develop a schema validator and an evidence-attribution checker that
independently diagnose structural and evidential failure in generated RAG output.

\textbf{RO3.} Develop a typed corrective-routing controller that maps each diagnosed
failure condition to a distinct action: local regeneration, entity-constrained
re-retrieval, or abstention.

\textbf{RO4.} Implement and evaluate SAFE-RAG against five baseline conditions and six
metrics on the ObliQA and FinDER/FinanceBench corpora, with every condition run under a
matched correction budget, and test pre-registered hypotheses H1-H3.

% ================================================================ 7
\section{Scope of Research}

The project is scoped to financial and regulatory compliance document question answering.
The primary evaluation corpus is ObliQA [15], comprising 27,869 question-and-passage pairs
drawn from real Abu Dhabi Global Markets financial regulation, supplemented by FinDER [16]
(5,703 expert-annotated analyst-style query, evidence, and answer triplets drawn from
S\&P 500 filings) and by filings from the FinanceBench benchmark [17]. A subset of both
corpora will be extended into schema-bound question-and-answer pairs, and a smaller
adversarial test set will be hand-constructed to contain deceptive-grounding-style
failures, which neither corpus contains by design.

The scope explicitly excludes construction of a new base corpus, and excludes empirical
demonstration of domain generality beyond finance; the latter is identified as a
direction for future work rather than a claim of the present project.

% ================================================================ 8
\section{Research Methodology}

\subsection{Systematic Review}

The proposal is grounded in a scoping review already conducted (May-July 2026) using
targeted Boolean searches across arXiv, Google Scholar, ACL Anthology, and IEEE Xplore,
from which 30 papers meeting defined inclusion criteria were selected for full-text
review. A supplementary update sweep was run in July 2026 against arXiv listings for
cs.CL and cs.IR, using the additional query terms (``failure-aware'' OR ``diagnose'' OR
``repair'') AND (``retrieval-augmented generation'' OR ``agentic RAG''), specifically to
capture typed-routing work published into the RAG setting during the review period. That
sweep returned five further papers meeting the same inclusion criteria [19], [20], [21],
[22], [23], bringing the total to 35. Three of the five materially narrow the gap claimed
by an earlier draft of this proposal, and their inclusion is reported explicitly here
rather than absorbed silently, because the publication rate in this area is itself a
threat to the validity of any gap claim (Section 8.4).

The funded project will formalize this scoping work into a complete, reproducible
systematic review following the PRISMA 2020 reporting standard and the Kitchenham and
Charters protocol for software-engineering literature reviews, including full
per-database screening counts and a supervisor-verified second screening pass to address
the single-reviewer limitation of the scoping stage.

\subsection{SAFE-RAG Architecture}

SAFE-RAG extends a standard retrieve-then-generate pipeline with two independently
diagnosed checks and a routing function. A hybrid retriever combines dense retrieval
scores in the style of Dense Passage Retrieval [13] with sparse lexical scores in the
style of BM25 [14], optionally followed by cross-encoder reranking, to return the top-k
passages for a query. A generator then produces a structured answer conditioned on the
query, the retrieved context, and a target schema.

A schema validator computes a binary pass/fail indicator of whether the generated answer
parses against the schema. On failure, the corrective action is local regeneration
against the same retrieved context, since the fault concerns output structure rather than
evidence quality. On schema pass, an evidence-attribution checker computes a composite
score combining an entailment component, following established faithfulness-scoring
approaches [7], [8], with an attribution component that verifies the retrieved evidence
concerns the same entity, clause, or category named in the query, following the
entity-attribution verification approach used to detect deceptive grounding [10]. A
routing function uses both checks, together with the number of correction attempts
already made, to decide whether to return the answer, regenerate, re-retrieve using an
entity-constrained reformulated query, or abstain after a configured maximum number of
attempts.

Two considerations justify diagnosing the conditions separately rather than collapsing
them into a single quality score. First, self-correction driven by no external, specific
signal of what is wrong frequently fails to improve output and can degrade it [12].
Second, repairability depends on failure type: format failures can be repaired by
rewriting over evidence that is already sufficient, whereas evidence failures cannot, an
asymmetry reported directly in the RAG repair literature [19]. A two-condition router is
the mechanism by which that asymmetry is exploited.

\subsection{Evaluation Design}

Five conditions will be compared: naive RAG with no schema enforcement and no
faithfulness check (a floor); schema-constrained generation only; faithfulness-checking
only; an undifferentiated single-retry baseline in which any detected failure triggers
the same generic retry regardless of type, using the same underlying validator and
checker components (the key comparison for RQ2); and the full SAFE-RAG system.

All five conditions will be run under a matched correction budget, defined as an equal
maximum number of generator and retriever calls per query. This control is not optional:
a repair loop compared against a single-pass baseline without budget control measures
additional compute rather than the value of the diagnosis [21], and results will not be
reported in a form that confounds the two. A secondary sweep will vary the maximum number
of correction attempts across one, two, and three to characterize how the advantage of
typing scales with the budget available.

Six metrics will be reported: schema validity rate; faithfulness score (entailment-based,
consistent with RePASs and AlignScore [8]); entity-attribution accuracy, a metric
introduced for this project; end-task correctness; repair efficiency (additional model or
retrieval calls required per query, within the matched budget); and abstention
precision/recall. To avoid the circularity of using a single automated judge both to
trigger repair and to score the result, the entailment model used inside the routing loop
and the model used to compute the reported faithfulness metric will be drawn from
different model families, and both will be validated against a human-annotated subsample
of at least 150 items before the main experiments are run.

Three pre-registered hypotheses will be tested:

\begin{itemize}
\item \textbf{H1.} On the adversarial deceptive-grounding test subset, SAFE-RAG achieves
entity-attribution accuracy at least 15 percentage points higher than the
faithfulness-checked-only baseline.

\item \textbf{H2.} Under a matched correction budget, SAFE-RAG achieves higher end-task
correctness than the undifferentiated single-retry baseline while consuming no more
additional model or retrieval calls per query on average. The budget-matching condition is
what makes this hypothesis falsifiable; without it, any observed difference is
attributable to additional compute [21].

\item \textbf{H3.} Ablating the typed-routing step, so any detected failure triggers the
same generic retry action using the same validator and checker components, reduces
end-task correctness relative to the full SAFE-RAG system, mirroring the ablation result
reported for typed routing in its original tool-calling domain [11] and its replication
within RAG [19].
\end{itemize}

\subsection{Limitations and Risk Mitigation}

Two risks specific to this proposal are recorded here, alongside the review limitations
noted in Section 8.1, so that their mitigations form part of the funded work plan.

The first concerns the durability of the gap claim itself. A contribution defined by an
absence in a fast-moving literature decays measurably over months: three of the five
papers added by the July 2026 update sweep were posted within the four months preceding
this submission, and two of them narrowed the gap claimed by an earlier draft of this
document. The mitigation is procedural rather than analytical. The claim is stated at the
narrowest defensible scope rather than the broadest available one, restricted to the
conjunction of schema invalidity and evidence misattribution in a compliance setting; a
standing arXiv alert on the four query clusters used in the review will be maintained for
the duration of the project; and the capability coverage assessment described in Section 4
will be re-run at each milestone review. A gap claim in this area should be read as a
dated observation, not as a stable property of the literature.

The second concerns the base rate of the failure mode the attribution branch targets.
Because the adversarial subset described in Section 7 is hand-constructed, its
distribution does not by itself establish how often deceptive grounding occurs in
practice. The rate at which it arises naturally in ObliQA and FinDER will therefore be
measured and reported as a descriptive statistic in its own right, and a low measured rate
would weaken the practical motivation for the attribution branch even if the H1 result
holds on the adversarial subset. That measurement is scheduled early, in Months 1 to 2,
precisely so that a negative result arrives before the implementation phase commits
resources, allowing the balance of effort between the schema-validity and attribution
branches to be revised while the project can still absorb the change.

% ================================================================ 9
\section{Expected Outcomes}

\begin{itemize}
\item A validated SAFE-RAG framework and an open research prototype comprising the schema
validator, evidence-attribution checker, and typed routing controller.

\item A completed PRISMA 2020-compliant systematic review formalizing the present
35-paper scoping review.

\item A reusable adversarial deceptive-grounding test set for financial and regulatory RAG
evaluation, together with a reported measurement of the natural rate of deceptive
grounding in ObliQA and FinDER.

\item A manuscript reporting the framework and evaluation results, targeted at a
regulatory natural language processing or applied machine learning venue, with submission
planned before the end of the project.

\item A direct contribution toward the applicant's graduate research program, extending
the applicant's prior schema-validated self-correction prototype into a formally evaluated
contribution.
\end{itemize}

% ================================================================ 10
\section{Research Contributions}

This project contributes: (1) a 35-paper scoping literature review scored against five
capability dimensions, establishing that the conjunction of schema-constrained output,
evidence-attribution checking, and typed corrective routing is currently unoccupied, while
recording explicitly that typed routing on its own is now established within RAG;
(2) SAFE-RAG, an architecture that performs schema validation and evidence-attribution
checking as two independently diagnosed conditions and routes each detected failure type
to a distinct corrective action; (3) an evaluation plan with pre-registered, falsifiable
hypotheses, run under a matched correction budget and built on existing regulatory and
financial benchmarks rather than a newly constructed corpus; and (4) a concrete deployment
pathway aligned with Malaysia's regulatory-technology development priorities [18].

% ================================================================ 11
\Needspace*{6\baselineskip}
\section{Project Timeline}

The timeline below is deliberately compressed to 12 months rather than the multi-year
horizon typical of a new framework proposal. This is possible because the scoping review,
gap analysis, and architecture design described in this proposal are already complete; the
funded period begins directly with implementation. The literature review, framework
design, and evaluation plan are treated as parallel, continuously updated work rather than
sequential gates, so that a submittable manuscript can be assembled from Month 9 onward
rather than only after every phase below has fully closed. The literature-monitoring task
described in Section 8.4 runs across the full period rather than concluding with the
Month 2 review.

\vspace{6pt}

{\setlength{\parskip}{0pt}
\renewcommand{\arraystretch}{1.25}
\setlength{\tabcolsep}{5pt}
\fontsize{9.5}{12}\selectfont
\begin{tabular}{|>{\bfseries}m{90pt}|>{\centering\arraybackslash}m{60pt}|p{280pt}|}
\hline
\rowcolor{ruleblue}
\textcolor{white}{\fontsize{10}{12}\selectfont\textbf{Phase}} &
\textcolor{white}{\fontsize{10}{12}\selectfont\textbf{Months}} &
\textcolor{white}{\fontsize{10}{12}\selectfont\textbf{Milestone / Deliverable}} \\ \hline
1. Systematic review \& design & 1--2 & Full PRISMA 2020-compliant systematic review
completed (extending the 35-paper scoping review already conducted for this proposal);
schema design finalized; adversarial deceptive-grounding test subset constructed; natural
rate of deceptive grounding in ObliQA and FinDER measured and reported. \\ \hline
2. Core diagnostic components & 3--5 & Schema validator and evidence-attribution checker
implemented and unit-validated against a human-annotated subset; entailment models for the
routing loop and for the reported faithfulness metric selected from different model
families. \\ \hline
3. Typed routing controller & 6--8 & Rule-based routing policy implemented and integrated
with the diagnostic components; learned routing policy explored if time permits. \\ \hline
4. Evaluation & 9--10 & All five baseline conditions run on ObliQA and
FinDER/FinanceBench under a matched correction budget; hypotheses H1--H3 tested against
pre-registered criteria. \\ \hline
5. Analysis \& publication & 11--12 & Results analysed; manuscript prepared and submitted
to a regulatory NLP or applied machine learning venue. \\ \hline
\end{tabular}\par}

% ================================================================ 12
\section{Conclusion}

This proposal identifies a specific, evidenced gap in the current literature. Typed error
routing, established first for tool-calling agents and carried during 2026 into agentic
RAG trajectories, is a demonstrated mechanism rather than an open question, and this
proposal treats it as a premise. What no reviewed system does is diagnose output structure
and evidence attribution as two independently detectable conditions over the same
structured output and route each to a corrective action selected on the basis of which
condition failed. The existing diagnosis vocabularies, built for trajectory-level repair,
cannot express a faithful-but-misattributed answer at all, and no typed-routing system in
this line has been evaluated on financial or regulatory content. SAFE-RAG addresses that
gap directly, diagnosing schema invalidity and evidence misattribution as two independently
detected conditions and routing each to a distinct corrective action, while building on
existing regulatory and financial benchmarks rather than requiring new data collection.

The evaluation plan specifies falsifiable, pre-registered hypotheses against both an
undifferentiated-retry baseline and a faithfulness-only baseline, under a correction budget
matched across all conditions, directly testing whether typed diagnosis over this
particular pair of conditions, rather than correction activity or additional compute in
general, accounts for any observed improvement. Within a compressed 12-month timeline, the
project is designed to extend the applicant's prior schema-validated self-correction
prototype into a formally evaluated, publishable research contribution, with an initial
deployment context aligned to Malaysia's regulatory-technology priorities and a technical
contribution intended to generalize to other high-stakes, structured-output domains in
future work.

% ================================================================ references
\section*{References}

{\setlength{\parskip}{6pt}\fontsize{9.5}{12}\selectfont

\refitem{1}{Z. Chen et al., ``FinQA: A Dataset of Numerical Reasoning over Financial
Data,'' in \textit{Proc. EMNLP}, 2021.}
\refitem{2}{F. Zhu et al., ``TAT-QA: A Question Answering Benchmark on a Hybrid of
Tabular and Textual Content in Finance,'' in \textit{Proc. ACL}, 2021.}
\refitem{3}{Z. Chen et al., ``ConvFinQA: Exploring the Chain of Numerical Reasoning in
Conversational Finance Question Answering,'' in \textit{Proc. EMNLP}, 2022.}
\refitem{4}{Y. Zhao et al., ``MultiHiertt: Numerical Reasoning over Multi Hierarchical
Tabular and Textual Data,'' in \textit{Proc. ACL}, 2022.}
\refitem{5}{S. Wu et al., ``BloombergGPT: A Large Language Model for Finance,''
arXiv:2303.17564, 2023.}
\refitem{6}{``XGrammar: Flexible and Efficient Structured Generation Engine for Large
Language Models,'' arXiv:2411.15100, 2024.}
\refitem{7}{S. Es, J. James, L. Espinosa-Anke, and S. Schockaert, ``RAGAS: Automated
Evaluation of Retrieval Augmented Generation,'' arXiv:2309.15217, 2023.}
\refitem{8}{Y. Zha, Y. Yang, R. Li, and Z. Hu, ``AlignScore: Evaluating Factual
Consistency with a Unified Alignment Function,'' in \textit{Proc. ACL}, 2023.}
\refitem{9}{``Faithfulness-Aware Uncertainty Quantification for Fact-Checking
Retrieval-Augmented Generation (FRANQ),'' arXiv:2505.21072, 2026.}
\refitem{10}{``Deceptive Grounding: Entity Attribution Failure in Clinical
Retrieval-Augmented Generation,'' arXiv:2607.09349, 2026.}
\refitem{11}{``PALADIN: Self-Correcting Language Model Agents to Cure Tool-Failure
Cases,'' in \textit{Proc. ICLR}, arXiv:2509.25238, 2026.}
\refitem{12}{``Large Language Models Cannot Self-Correct Reasoning Yet,'' in
\textit{Proc. ICLR}, arXiv:2310.01798, 2024.}
\refitem{13}{V. Karpukhin et al., ``Dense Passage Retrieval for Open-Domain Question
Answering,'' in \textit{Proc. EMNLP}, 2020.}
\refitem{14}{S. Robertson and H. Zaragoza, ``The Probabilistic Relevance Framework: BM25
and Beyond,'' Foundations and Trends in Information Retrieval, vol. 3, no. 4,
pp. 333-389, 2009.}
\refitem{15}{``RIRAG: Regulatory Information Retrieval and Answer Generation,'' in
\textit{Proc. EMNLP}, arXiv:2409.05677, 2024.}
\refitem{16}{``FinDER: A Financial Dataset for Question Answering and Evaluating
Retrieval-Augmented Generation,'' in \textit{Proc. ICLR Workshop}, arXiv:2504.15800, 2025.}
\refitem{17}{P. Islam et al., ``FinanceBench: A New Benchmark for Financial Question
Answering,'' arXiv:2311.11944, 2023.}
\refitem{18}{Bank Negara Malaysia, ``Financial Sector Blueprint 2022 to 2026,'' Kuala
Lumpur, Malaysia, 2022.}
\refitem{19}{``Doctor-RAG: A Failure-Aware Repair Framework for Agentic
Retrieval-Augmented Generation,'' arXiv:2604.00865, 2026.}
\refitem{20}{``Skill-RAG: Failure-State-Aware Retrieval Augmentation via Hidden-State
Probing and Skill Routing,'' arXiv:2604.15771, 2026.}
\refitem{21}{``Diagnosing and Repairing Factual Errors in RAG under Budget Constraints,''
arXiv:2606.29377, 2026.}
\refitem{22}{K. K. Leung, M. Belbahri, Y. Sui, A. Labach, X. Zhang, S. A. Rose, and
J. C. Cresswell, ``Classifying and Addressing the Diversity of Errors in
Retrieval-Augmented Generation Systems,'' in \textit{Proc. EACL}, arXiv:2510.13975, 2026.}
\refitem{23}{``Citation-Closure Retrieval and Per-Rule Attribution for Real-World
Regulatory Compliance Question Answering,'' arXiv:2605.29742, 2026.}

}

\end{document}
